%% kapa: exam=2 round=26 year=2015 period=2 subject=appraisal_theory subject_ko=감정평가이론 seq=2 number=2 points=30
%% source: Q-Net 감정평가사 기출문제 (cst00309) artlSeq=5149601 file=extracted/1008/5149601/2015년 제26회 감정평가사 2차시험 기출문제/2교시(감정평가이론).pdf
%% conversion: kordoc 4.24.51 -> tools/gen2.py (auto); review=proofread-gpt-sol; flags=none
\begin{longquestion}{【문제 2】}
감정평가목적에 따라 감정평가금액의 격차가 큰 경우가 있다. 다음 물음에 답하시오. (30점) \par
1) 보상평가, 경매평가, 담보평가의 목적별 평가방법을 약술하고, 동일한 물건이 감정평가목적에 따라 감정평가금액의 격차가 큰 사례 5가지를 제시하고 그 이유를 설명하시오. (20점) \par
2) 주거용 건물을 신축하기 위해 건축허가를 득하여 도로를 개설하고 입목을 벌채중인 임야를 평가하고자 한다. 개발 중인 토지의 평가방식에는 공제방식과 가산방식이 있다. 공제방식은 개발 후 대지가격에서 개발에 소요되는 제반비용을 공제하는 방식이고, 가산방식은 소지가격에 개발에 소요되는 비용을 가산하여 평가하는 방식이다. 두 가지 방식에 따른 감정평가금액의 격차가 클 경우 보상평가, 경매평가, 담보평가에서 각각 어떻게 평가하는 것이 더 적절한지 설명하시오. (10점) \par
\end{longquestion}
